% Primary and supplementary source basis: :contentReference[oaicite:14]{index=14} :contentReference[oaicite:15]{index=15} :contentReference[oaicite:16]{index=16}

\section{Cauchy--Riemann Equations}

\begin{conceptbox}
    Let
    \[
        f(z)=u(x,y)+iv(x,y), \qquad z=x+iy.
    \]
    If $f$ is differentiable at a point, then
    \[
        \frac{\partial u}{\partial x}=\frac{\partial v}{\partial y},
        \qquad
        \frac{\partial u}{\partial y}=-\frac{\partial v}{\partial x}.
    \]
    These are the \textbf{Cauchy--Riemann equations}.
\end{conceptbox}

\subsection*{Motivation / Intuition}

Because the derivative in the complex plane must be independent of direction, differentiating in the $x$-direction and in the $y$-direction must give the same result. That compatibility condition is exactly what produces the Cauchy--Riemann equations.

\subsection*{Derivation}

Starting from
\[
    f'(z)=\lim_{\Delta z\to 0}\frac{\Delta w}{\Delta z},
\]
take $\Delta z=\Delta x$ first. Then
\[
    f'(z)=\frac{\partial u}{\partial x}+i\frac{\partial v}{\partial x}.
\]

Now take $\Delta z=i\Delta y$. Then
\[
    f'(z)=\frac{\partial v}{\partial y}-i\frac{\partial u}{\partial y}.
\]

Since both expressions must be equal,
\[
    \frac{\partial u}{\partial x}+i\frac{\partial v}{\partial x}
    =
    \frac{\partial v}{\partial y}-i\frac{\partial u}{\partial y}.
\]

Equating real and imaginary parts gives
\[
    \frac{\partial u}{\partial x}=\frac{\partial v}{\partial y},
    \qquad
    \frac{\partial v}{\partial x}=-\frac{\partial u}{\partial y}.
\]

\begin{conceptbox}
    If the partial derivatives are continuous in a neighborhood and the Cauchy--Riemann equations hold there, then $f$ is differentiable there.
\end{conceptbox}

\subsection*{Polar Form}

If
\[
    f(z)=u(r,\theta)+iv(r,\theta),
\]
then the Cauchy--Riemann equations in polar coordinates are
\[
    \frac{\partial u}{\partial r}
    =
    \frac{1}{r}\frac{\partial v}{\partial \theta},
    \qquad
    \frac{\partial v}{\partial r}
    =
    -\frac{1}{r}\frac{\partial u}{\partial \theta}.
\]

These are useful when the function is naturally expressed in polar form.

\subsection*{Worked Examples}

\begin{examplebox}
    \textbf{Example 1: Verify analyticity of $f(z)=z^2$}

    Write
    \[
        z^2=(x+iy)^2=(x^2-y^2)+i(2xy).
    \]

    So
    \[
        u(x,y)=x^2-y^2,
        \qquad
        v(x,y)=2xy.
    \]

    Now compute the partial derivatives:
    \[
        u_x=2x,
        \qquad
        u_y=-2y,
        \qquad
        v_x=2y,
        \qquad
        v_y=2x.
    \]

    Check the Cauchy--Riemann equations:
    \[
        u_x=v_y \quad \Rightarrow \quad 2x=2x,
    \]
    \[
        u_y=-v_x \quad \Rightarrow \quad -2y=-2y.
    \]

    They hold everywhere, and the partial derivatives are continuous everywhere. Therefore, $z^2$ is entire.

    The derivative is
    \[
        f'(z)=u_x+iv_x=2x+i2y=2z.
    \]
\end{examplebox}

\begin{examplebox}
    \textbf{Example 2: Test $f(z)=|z|^2$}

    Write
    \[
        |z|^2=x^2+y^2+i0.
    \]

    Thus
    \[
        u=x^2+y^2,
        \qquad
        v=0.
    \]

    Compute:
    \[
        u_x=2x,
        \qquad
        u_y=2y,
        \qquad
        v_x=0,
        \qquad
        v_y=0.
    \]

    The equations require
    \[
        2x=0,
        \qquad
        2y=0.
    \]

    So they hold only at $(0,0)$. Hence $|z|^2$ is not analytic in any neighborhood of any point, and therefore is not analytic anywhere.
\end{examplebox}

\begin{examplebox}
    \textbf{Example 3: Use the polar form for $f(z)=1/z$}

    Write
    \[
        f(z)=\frac{1}{z}=\frac{1}{r}e^{-i\theta}
        =
        \frac{\cos\theta}{r}
        -i\frac{\sin\theta}{r}.
    \]

    So
    \[
        u(r,\theta)=\frac{\cos\theta}{r},
        \qquad
        v(r,\theta)=-\frac{\sin\theta}{r}.
    \]

    Compute:
    \[
        u_r=-\frac{\cos\theta}{r^2},
        \qquad
        v_\theta=-\frac{\cos\theta}{r},
    \]
    hence
    \[
        \frac{1}{r}v_\theta
        =
        -\frac{\cos\theta}{r^2}
        =
        u_r.
    \]

    Also,
    \[
        v_r=\frac{\sin\theta}{r^2},
        \qquad
        u_\theta=-\frac{\sin\theta}{r},
    \]
    so
    \[
        -\frac{1}{r}u_\theta
        =
        \frac{\sin\theta}{r^2}
        =
        v_r.
    \]

    Thus the polar Cauchy--Riemann equations hold for $r\neq 0$. Therefore, $1/z$ is analytic for $z\neq 0$.
\end{examplebox}

\begin{noteBox}
    The Cauchy--Riemann equations are among the most important tests in engineering complex analysis. They help detect whether a candidate field can behave like a valid complex potential, analytic transfer function, or conformal mapping.
\end{noteBox}

\subsection*{Common Mistakes / Notes}

\begin{itemize}
    \item Satisfying the Cauchy--Riemann equations at a single point does not guarantee analyticity there.
    \item Continuity of the partial derivatives in a neighborhood is the standard practical condition used with the equations.
    \item The equations test complex differentiability, not ordinary differentiability in $\mathbb{R}^2$.
\end{itemize}